\section{Framework: MDI and FIP}
\label{sec:framework}

\paragraph{(i) A judge score is a draw from a distribution.}
Fix an evaluation environment
\begin{equation*}
  \env = (\text{judge model},\ \text{prompt template},\ \text{task},\
  \text{scale format},\ \text{temperature})
\end{equation*}
and an output $x$. The judge's score is not a constant of $(x, \env)$. Repeat
the same call $r = 1, \dots, N$ times and it returns draws $s^{(r)}(x;\env)$
from a distribution. Everything below follows from keeping that distribution
instead of using its first value. Scores are reported scale-normalized (\%p;
raw values in the appendix). We write $\bar{s}_N$ for a system's item-mean
score over $N$ repeats and $\Delta = \bar{s}_N^A - \bar{s}_N^B$ for the
observed gain.

\paragraph{(ii) $\sigma(\env)$ is observable only through same-input repeats.}
The spread $\sigma(\env)$ cannot be read off a single run, a leaderboard or a
model card: it is observable only by scoring the \emph{same input} repeatedly
in the \emph{same environment}. That is the one idea we take from psychometrics
(test--retest reliability), and Section~\ref{sec:experiments} is built to
produce such repeats.

\paragraph{(iii) The prediction interval of an observed gain.}
When the true gain is zero, the observed $\Delta = \bar{s}_N^A - \bar{s}_N^B$
still has a distribution, inherited from the judge's. Inside its central 95\%
prediction interval, a gain cannot be told apart from the judge scoring the
same thing again. That interval depends on $(\env, N)$ and not on which
systems are compared. Definition~\ref{def:mdi} states this as the null
distribution $D_0(\env, N)$, estimated nonparametrically from
repeated-scoring data.

\paragraph{(iv) $N$-dependence and the noise floor.}
The interval shrinks with the measurement budget $N$, but not to zero. The
variance of a mean of $N$ repeats decomposes as
\begin{equation}
  \label{eq:decay}
  \operatorname{Var}\!\left[\bar{s}_N\right]
    \;=\; \frac{\sigma^2(\env)}{N} \;+\; \omega^2(\env),
\end{equation}
the standard variance-components form (cf.\
\citealp{miller2024errorbars, messing2026tee}). Repeats reduce
$\sigma^2(\env)$; they do not touch the floor $\omega^2(\env)$. Two things
follow: the threshold depends on $N$, and buying more repeats stops paying at
$\omega(\env)$. Experiment~3's sweep fits Eq.~\eqref{eq:decay} per environment
to estimate both, and its $\omega^2$ independently replicates the CI
half-width floor of \citet{messing2026tee}.

\paragraph{(v) Everything depends on the environment.}
Every quantity above depends on $\env$. A threshold computed for one judge,
prompt, task, and scale does not carry over to another; the grid of
Section~\ref{sec:experiments} measures how badly it fails to carry over, which
is why every table is indexed by environment.
An environment includes the scoring harness (parsing, rubric rendering,
retries, provider), so in measurement-theory terms our intrinsic axis is
\emph{repeatability} and our procedural axis \emph{reproducibility}.

\paragraph{(vi) A procedure, not a benchmark.}
A fixed-item benchmark would contradict the thesis that every quantity here
depends on the environment. What we offer is the procedure a team runs to get
its own numbers, plus the implementation; the tables below illustrate it rather
than replace it.

\subsection{FIP and MDI}
\label{sec:framework-defs}

\begin{definition}[False Improvement Probability]
  \label{def:fip}
  For an observed gain $\Delta$ measured in environment $\env$ with $N$
  scoring repeats per system,
  \begin{equation*}
    \fip(\Delta, \env, N)
      \;=\; P\bigl(\Delta' \le 0 \,\big|\, \text{observed } \Delta,\ \env,\ N\bigr),
  \end{equation*}
  where $\Delta'$ is the gain re-measured under an independent replication
  in the same environment with the same budget. $\fip$ is estimated
  nonparametrically from repeated-scoring data (Experiment~2).
\end{definition}

Conditioning on the observed effect is the move \citet{baumann2025llmhacking}
make for annotation-stage $p$-values; we make it for measurement-stage
score gains. The ancestor of both is Card et al.'s Type-S error
\citep{card2020power}, with the stochasticity relocated from the sample to
the scorer.

\begin{definition}[Minimum Detectable Improvement]
  \label{def:mdi}
  Let $D_0(\env, N)$ be the distribution of mean score differences between
  two independent $N$-repeat scorings of the \emph{same} systems in
  environment $\env$ (pairs whose true skill difference is zero by
  construction). Then
  \begin{equation*}
    \mdi(\env, N)
      \;=\; \tfrac{1}{2}\,
      \operatorname{width}\!\bigl(\mathrm{PI}_{95\%}\bigl[D_0(\env, N)\bigr]\bigr),
  \end{equation*}
  the half-width of the central 95\% prediction interval of the null
  distribution. Via Eq.~\eqref{eq:decay}, $\mdi$ falls with $N$ at a rate
  set by $\sigma^2(\env)$ and flattens onto the floor $\omega(\env)$: the
  dependence on $N$ is part of the definition.
\end{definition}

$D_0$ is built by splitting each cell's estimation repeats, per item, into two
disjoint pools and collecting the difference of their $N$-repeat means
(Appendix~\ref{app:extra}).

We \emph{define} rather than derive $\mdi$, for a practical reason: the null
interval is exactly the region in which the instrument cannot tell two
identical systems apart. One property should be stated plainly: a true gain
of exactly $\mdi$ is measured above $\mdi$ only about half the time. So $\mdi$
controls false claims at the stated level $\alpha$, not the chance of catching a real gain.
A team that wants an $80\%$ chance needs a larger gain, and we measure the
factor rather than assume it: shifting $D_0$ and reading where $80\%$ of it
clears $\mdi$ gives $1.43\,\mdi$ (median over 36 cells, range
$1.34$--$1.46$; the closed-form cross-check is
Appendix~\ref{app:extra}). We tabulate the $\alpha$-level value
because $D_0$ measures it directly, and carry the $80\%$ column beside it.
The FIP-crossing read-off
$\min\{\Delta : \fip(\Delta, \env, N) \le \alpha\}$ is a consistency probe,
not a second definition; it sits \emph{above} the null-PI value and
Section~\ref{sec:experiments} reports the gap. Every decision here uses the
null-PI $\mdi$.

The nearest prior statistics describe the instrument, not the decision. Feuer
et al.'s $\mathrm{Sens}$ \citep{feuer2025judgmentnoise} carries no
$N$-dependence, and $\Delta^*_{75}$ of \citet{usami2026darkcurrent} is measured
per response pair from single calls: $\Delta^*(J)$, where $\mdi$ is $\Delta^*(J, N)$ at
a stated error level. The MDE lineage \citep{card2020power,
miller2024errorbars, kotawala2026resolution} places the variance in the sample,
ours in the scorer.

